\section{System Model}
\label{sec:system}

As illustrated in \cref{fig:system}, we consider a post-disaster network in which a single UAV complements a damaged ground relay network in forwarding alerts to a rescue center. \Cref{tab:notation} summarizes the main notation.

\begin{figure}[htbp]
\centering
\resizebox{0.9\textwidth}{!}{% System-model scene (Figure 1), drawn with TikZ and fontawesome5 icons.
\begin{tikzpicture}[x=1cm, y=1cm, >={Stealth[length=2.2mm]},
    lbl/.style={font=\footnotesize, align=center},
    tag/.style={font=\scriptsize, align=center, fill=groundgreen, inner sep=1pt},
    backhaul/.style={draw=black!55, line width=1.4pt},
    failed/.style={draw=red!75!black, line width=1.4pt, dash pattern=on 4pt off 3pt},
    uplink/.style={->, draw=red!80!black, line width=1.1pt},
    downlink/.style={->, draw=blue!70!black, line width=1.1pt, dash pattern=on 7pt off 4pt},
    ground/.style={->, draw=green!45!black, line width=1.1pt},
    traj/.style={draw=black, line width=0.7pt, dash pattern=on 5pt off 2pt on 1pt off 2pt}]
  % ---------------------------------------------------------------- ground
  \fill[groundgreen] (6.6,3.4) ellipse (6.6 and 3.1);
  \draw[black!45] (6.6,3.4) ellipse (6.6 and 3.1);
  % damage zone
  \fill[red!55!brown, opacity=0.22] (9.6,3.3) ellipse (2.7 and 1.75);
  \draw[red!55!black, dashed, line width=0.6pt] (9.6,3.3) ellipse (2.7 and 1.75);
  \node[lbl, text=red!50!black, fill=white, fill opacity=0.75, text opacity=1, rounded corners=2pt, inner sep=1.5pt] at (10.4,1.35) {Damage zone};
  % ---------------------------------------------------------------- ground relay network
  \coordinate (C)  at (1.6,3.3);
  \coordinate (T1) at (3.4,5.0);
  \coordinate (T2) at (3.6,1.9);
  \coordinate (T3) at (5.6,3.6);
  \coordinate (T4) at (7.0,1.5);
  \coordinate (T5) at (7.6,5.0);
  \draw[backhaul] (C) -- (T1) (C) -- (T2) (T1) -- (T3) (T2) -- (T3) (T2) -- (T4) (T3) -- (T5) (T3) -- (T4);
  \draw[failed] (T5) -- (9.0,4.3);
  \node[text=red!75!black, font=\large] at (8.25,4.68) {\faTimes};
  \foreach \t in {T1,T2,T3,T4,T5} \node[text=black!70, font=\LARGE] at (\t) {\faBroadcastTower};
  \node[text=black!70, font=\LARGE] at (9.0,4.3) {\faBroadcastTower};
  \node[tag, text=red!60!black, fill=none] at (8.95,3.82) {lost backhaul};
  \node[tag] at (3.4,5.55) {relay};
  \node[tag] at (7.0,0.95) {relay $v$};
  \node[tag] at (5.0,3.2) {drop node $u$};
  % rescue center
  \node[text=red!70!black, font=\Huge] at (C) {\faHospital};
  \node[lbl, anchor=north] at (1.6,2.65) {Rescue center $c$};
  % ---------------------------------------------------------------- buildings in the zone
  \node[text=brown!70!black, font=\huge] at (8.2,2.4) {\faHouseDamage};
  \node[text=brown!70!black, font=\huge] at (10.2,4.4) {\faHouseDamage};
  \node[text=brown!70!black, font=\huge] at (11.3,2.6) {\faHouseDamage};
  \node[text=black!50, font=\huge] at (4.9,5.3) {\faBuilding};
  \node[text=black!50, font=\huge] at (5.2,0.9) {\faCity};
  % ---------------------------------------------------------------- sources raising alerts
  \foreach \x/\y/\n in {9.2/3.1/s1, 10.3/3.2/s2, 9.7/2.2/s3, 11.6/3.5/s4} {
    \node[text=blue!55!black, font=\LARGE] (\n) at (\x,\y) {\faUsers};
    \node[text=orange!90!black, font=\small] at (\x+0.38,\y+0.38) {\faExclamationTriangle};
  }
  \node[lbl] at (9.7,1.72) {sources};
  % ---------------------------------------------------------------- UAV
  \begin{scope}[shift={(8.7,7.9)}, scale=0.62]
    \fill[black!80] (-1.1,0.05) rectangle (1.1,0.22);
    \foreach \x in {-1.1,1.1} {
      \fill[black!80] (\x,0.13) circle (0.16);
      \fill[black!65] (\x,0.42) ellipse (0.62 and 0.11);
      \fill[black!80] (\x-0.04,0.13) rectangle (\x+0.04,0.42);
    }
    \fill[black!85] (0,0.0) ellipse (0.62 and 0.36);
    \fill[black!60] (0,-0.05) ellipse (0.28 and 0.16);
    \foreach \r in {0.45,0.7,0.95} \draw[black!80, line width=1.1pt] ($(0,-0.5)+(200:\r)$) arc (200:340:\r);
  \end{scope}
  \node[lbl] at (10.75,8.05) {UAV relay};
  \coordinate (U) at (8.7,7.35);
  % ---------------------------------------------------------------- trajectory
  \draw[traj, ->] (1.6,3.9) .. controls (2.0,6.6) and (5.2,7.6) .. (U);
  \draw[traj, ->] (U) to[out=0, in=100] (12.3,5.2) to[out=-80, in=0] (8.0,0.25) to[out=180, in=-90] (1.45,2.72);
  \node[lbl, rotate=28] at (3.5,6.75) {take-off};
  \node[text=red!75!black, font=\Large] at (12.3,5.5) {\faMapMarker*};
  \node[lbl, anchor=east] at (12.05,5.85) {hover stop};
  % ---------------------------------------------------------------- links
  \draw[uplink] (s1.north) -- ($(U)+(-0.15,0)$);
  \draw[uplink] (s4.north) -- ($(U)+(0.25,0)$);
  \draw[downlink] ($(U)+(-0.35,0)$) -- ($(T3)+(0.3,0.35)$);
  \draw[ground] (s3.west) -- ($(T4)+(0.35,0.1)$);
  % elevation angle at a source
  \draw[black!80, densely dotted] (11.6,3.78) -- (10.55,3.78);
  \draw[black!80] ($(11.6,3.78)+(180:0.75)$) arc (180:126:0.75);
  \node[font=\small] at ($(11.6,3.78)+(153:1.05)$) {$\theta$};
  % ---------------------------------------------------------------- legend
  \begin{scope}[shift={(13.2,8.4)}]
    \node[lbl, anchor=west] at (0,0) {Trajectory};
    \draw[traj, ->] (0,-0.35) -- (1.7,-0.35);
    \node[lbl, anchor=west] at (0,-0.85) {Uplink (access)};
    \draw[uplink] (0,-1.2) -- (1.7,-1.2);
    \node[lbl, anchor=west] at (0,-1.7) {Downlink};
    \draw[downlink] (0,-2.05) -- (1.7,-2.05);
    \node[lbl, anchor=west] at (0,-2.55) {Ground access};
    \draw[ground] (0,-2.9) -- (1.7,-2.9);
    \node[lbl, anchor=west] at (0,-3.4) {Backhaul / failed};
    \draw[backhaul] (0,-3.75) -- (0.75,-3.75);
    \draw[failed] (0.95,-3.75) -- (1.7,-3.75);
  \end{scope}
\end{tikzpicture}
}
\caption{Model of the UAV-assisted alert delivery system.}
\label{fig:system}
\end{figure}

\subsection{Network and Alert Model}
\label{sec:model}

The post-disaster ground relay network is represented by a graph $G=(V,E)$ taken from Topology Zoo~\cite{knight2011zoo} and normalized into a square region of side $L_{\mathrm{box}}$ with the same scaling factor on both axes. Node $c\in V$ is the rescue center and all other nodes are relays. A surviving edge $e\in E$ has backhaul capacity $C_e$, and a failed edge has capacity $0$. Alerts are raised at a set of sources $S$, which do not belong to $G$ and can transmit only over radio links. Each alert $a\in A$ has a source $s_a\in S$, a release slot $r_a$, a deadline $d_a$, and a size of $L_a$ bits. The mission consists of $N$ slots of length $\Delta$, and each slot is split into an access phase of length $\tau\Delta$, in which sources transmit to relays or to the UAV, and a downlink phase of length $(1-\tau)\Delta$, in which the UAV transmits to nodes of $G$.

Each alert has at most $K$ candidate paths of two types: a ground path $s_a\to v\leadsto c$ through a relay $v$ in range, and a UAV path $s_a\to\text{UAV}\to u\leadsto c$ with drop node $u\in V$. The segment $v\leadsto c$ is a minimum-hop backhaul path over the surviving edges. The candidate set consists of at most two ground paths with the lowest estimated cost, complemented by the UAV paths with the lowest cost; it is fixed before planning and shared by the baselines, the bound of \cref{sec:analysis}, and the initial plan of the proposed search. A plan assigns each alert a path $\pi_a$ or leaves it unserved; the proposed search may assign any loop-free route of either type over surviving edges, including routes outside the candidate set.

\begin{table}[htbp]
\centering
\small
\caption{Main notation.}
\label{tab:notation}
\begin{tabularx}{\textwidth}{@{}>{\raggedright\arraybackslash}p{0.2\textwidth}X@{}}
\toprule
\textbf{Symbol} & \textbf{Description} \\
\midrule
$G=(V,E)$, $c$ & Ground relay network and rescue center \\
$C_e$ & Backhaul capacity of edge $e$ (zero if failed) \\
$S$, $A$ & Set of sources and set of alerts \\
$s_a$, $r_a$, $d_a$, $L_a$ & Source, release slot, deadline, and size (bits) of alert $a$ \\
$N$, $\Delta$, $\tau$ & Number of slots, slot length, and access-phase fraction \\
$q[n]$, $H$, $V_{\max}$ & UAV horizontal position, altitude, and maximum speed \\
$P_k^{L}[n]$ & LoS probability between the UAV and ground point $k$ in slot $n$ \\
$R(b,s)$, $B_{\mathrm{tot}}$ & Rate with bandwidth $b$ and SNR coefficient $s$; bandwidth per phase \\
$K$, $x_{a,p}$, $\pi_a$ & Maximum number of candidate paths; candidate selection and path of alert $a$ \\
$b_{e,n}$, $w_{e,n}$ & Bandwidth and bandwidth weight of radio link $e$ in slot $n$ \\
$\omega$, $M$, $M_d$ & Channel realization; numbers of evaluation and design realizations \\
$z_{a,\omega}$, $\hat P_{\mathrm{timely}}$ & Timely indicator of alert $a$ in realization $\omega$; timely-delivery ratio \\
$\mathrm{Conn}$ & Fraction of sources with a time-respecting path to $c$ \\
$E$ & Evaluation budget of the search \\
\bottomrule
\end{tabularx}
\end{table}

\subsection{UAV Mobility and Channel Model}
\label{sec:channel}

The UAV flies at a fixed altitude $H$ along a trajectory $q[0],\dots,q[N]\in\R^2$ that satisfies
\begin{equation}
\label{eq:mobility}
q[0]=q_s,\qquad q[N]=q_t,\qquad \norm{q[n+1]-q[n]}\le V_{\max}\Delta,\quad n=0,\dots,N-1,
\end{equation}
and the representative position of slot $n$ is the midpoint $\bar q[n]=\bigl(q[n]+q[n+1]\bigr)/2$.

The link between the UAV and a ground point $k$ at position $w_k$ follows the PrLoS model of~\cite{huang2025ddatsap}:
\begin{equation}
\label{eq:geometry}
\begin{aligned}
d_k[n]&=\sqrt{\norm{\bar q[n]-w_k}^2+H^2},\qquad
\theta_k[n]=\frac{180}{\pi}\operatorname{atan2}\bigl(H,\norm{\bar q[n]-w_k}\bigr),\\
P_k^{L}[n]&=\frac{1}{1+a\exp\bigl[-b(\theta_k[n]-a)\bigr]},
\end{aligned}
\end{equation}
with channel gain $g^{L}=\beta_0 d^{-\alpha_L}$ under line of sight and $g^{N}=\mu\beta_0 d^{-\alpha_N}$ otherwise. For transmit power $p$ and noise power spectral density $N_0$, let $s=pg/N_0$ (in Hz); the rate achieved with bandwidth $b$ is
\begin{equation}
\label{eq:rate}
R(b,s)=b\log_2\!\Bigl(1+\frac{s}{b}\Bigr),\qquad R(0,s)=0.
\end{equation}
In a channel realization $\omega$, the LoS/NLoS state of the link to point $k$ in slot $n$ is drawn independently with probability $P_k^{L}[n]$, evaluated on the trajectory under consideration, and all plans are compared on common random numbers. The link from a source to a ground relay has a deterministic gain $\beta_0 d^{-\alpha_G}$ and exists only if $R(B_{\mathrm{tot}},s)\ge R_{\min}$. Sources transmit with power $p_s$, and the UAV splits a total power $P_U$ equally among at most two active downlinks per slot, so the power budget does not grow with the problem size.
